\chapter*{Tóm tắt}
\label{summary}

Phát hiện bất thường chuỗi thời gian là một bài toán được ứng dụng trong nhiều lĩnh vực quan trọng \cite{10.1145/3637528.3671936} \cite{NEURIPS2023_7a53bf4e} \cite{10.1145/3746252.3760799}. Tuy nhiên, chưa có nhiều phương pháp kết hợp \textit{học biểu diễn đa nhiệm (multi-task representation learning)} và \textit{cơ chế bộ nhớ (memory bank)}. Bên cạnh đó, \textit{ước lượng độ bất định nhận thức (epistemic uncertainty estimation)} của mô hình và \textit{thích ứng mô hình trên môi trường trực tuyến (online adaptation)} cũng là các vấn đề chưa được quan tâm đủ nhiều. Trong báo cáo này, chúng tôi đề xuất 2 toán tử truy vấn tương ứng với một bộ nhớ liên tục và một bộ nhớ rời rạc nhằm tính toán được biểu diễn chuyên biệt cho mỗi tác vụ. Chúng tôi cũng đề xuất 2 phiên bản ngẫu nhiên tương ứng cho hai toán tử đó nhằm ước lượng được độ bất định nhận thức của mô hình. Thêm vào đó, chúng tôi đề xuất một khung thích ứng trực tuyến dựa trên hàm mất mát tương phản offline--online và cơ chế \textit{bộ đệm xác minh (verification buffer)} để phát hiện và thích ứng trên các cửa sổ được xem là chứa \textit{mẫu hình bình thường mới (new normality)}. Kết quả thực nghiệm cho thấy phương pháp đề xuất có tốt hơn các phương pháp đối sánh trên các độ đo VUS-PR và Affiliation F1-score trong cả pha tiền huấn luyện ngoại tuyến và thích ứng trực tuyến. Thông qua phân tích thực nghiệm, chúng tôi kết luận rằng mức tăng trong các độ đo hiệu quả phần lớn đến từ hai mô-đun toán tử truy vấn (trong pha ngoại tuyến) và cơ chế làm mượt (trong pha trực tuyến). Ngoài ra, hai phiên bản ngẫu nhiên của hai toán tử truy vấn cũng cho thấy kết quả khả quan trong ước lượng độ bất định.